% Generated draft: requires mathematical and source review.
\documentclass[11pt]{article}
\usepackage[margin=25mm]{geometry}
\usepackage{fontspec}
\setmainfont{Latin Modern Roman}
\usepackage{amsmath,amssymb,mathtools,unicode-math}
\setmathfont{Latin Modern Math}
\usepackage{enumitem,needspace,xurl,hyperref}
\setlength{\parindent}{0pt}
\setlength{\parskip}{5pt}
\setlength{\emergencystretch}{3em}
\newcommand{\N}{\mathbb{N}}
\newcommand{\Z}{\mathbb{Z}}
\newcommand{\Q}{\mathbb{Q}}
\newcommand{\R}{\mathbb{R}}
\newcommand{\C}{\mathbb{C}}
\newcommand{\F}{\mathbb{F}}
\newcommand{\A}{\mathbb{A}}
\newcommand{\PP}{\mathbb{P}}
\newcommand{\E}{\mathbb{E}}
\newcommand{\eps}{\varepsilon}
\newcommand{\dd}{\,\mathrm{d}}
\begin{document}
\section*{UnsolvedMath ID: 701587}
\textbf{Source problem code:} catalogue entry
\subsection{Definitions and mathematical statement}
All groups considered below are finite. A formation is a nonempty class \(\mathfrak F\) of finite groups, closed under isomorphisms, homomorphic images, and subdirect products. Explicitly, the last condition means that if \(N_1,\ldots,N_r\trianglelefteq G\), where \(r\geq 1\), satisfy \(\bigcap_{i=1}^r N_i=1\) and \(G/N_i\in\mathfrak F\) for every \(i\), then \(G\in\mathfrak F\). A formation is hereditary if every subgroup of every member belongs to the formation. It is soluble if each of its members is soluble, meaning that its derived series terminates at the trivial subgroup.

Let \(\mathfrak U\) denote the formation of all supersoluble finite groups: these are the finite groups admitting a series
\[
1=G_0\trianglelefteq G_1\trianglelefteq\cdots\trianglelefteq G_n=G
\]
in which every \(G_i\) is normal in \(G\) and every factor \(G_i/G_{i-1}\) is cyclic of prime order. The trivial group is included by allowing \(n=0\).

For a class \(\mathfrak X\) of finite groups, a finite group \(G\) is \(\mathfrak X\)-critical if
\[
G\notin\mathfrak X
\quad\text{and}\quad
H\in\mathfrak X\ \text{for every proper subgroup }H<G.
\]
Write \(\Phi(G)\) for the Frattini subgroup, the intersection of all maximal proper subgroups of \(G\), with \(\Phi(1)=1\). A formation \(\mathfrak F\) is saturated if, for every finite group \(G\),
\[
G/\Phi(G)\in\mathfrak F\ \Longrightarrow\ G\in\mathfrak F.
\]

\textbf{Question.} Let \(\mathfrak F\) be any soluble hereditary formation of finite groups such that every finite \(\mathfrak F\)-critical group either has prime order or is \(\mathfrak U\)-critical. Must \(\mathfrak F\) be saturated? Equivalently, does this hypothesis imply
\[
\forall\,G\text{ finite},\qquad
G/\Phi(G)\in\mathfrak F\ \Longrightarrow\ G\in\mathfrak F?
\]
\subsection{Short English statement}
Suppose a subgroup-closed formation consists entirely of soluble finite groups, and every finite group minimally outside it is either of prime order or minimally outside the supersoluble groups. Must membership of a finite group's Frattini quotient force membership of the group itself?
\subsection{Sources}
\begin{itemize}
\item Local UnsolvedMath registry, permanent ID \(701587\), exact source code \texttt{catalogue entry}, title \emph{Kourovka Notebook Problem 19.39}. Supplied upstream revision: \texttt{d371ee917ac9647fbb06cd6fc529deec6f3dda82}.
\item S. F. Kamornikov, proposed in 2018, Problem 19.39 in \emph{The Kourovka Notebook (Unsolved Problems in Group Theory)}, 19th issue (2018). The supplied background also identifies the 21st issue edition (2026), arXiv:1401.0300v48. The supplied collection description names E. I. Khukhro and V. D. Mazurov as editors.
\item The local record labels the problem \emph{open} and records an update date of 2026-10-16. This is supplied registry information, not an independently verified assessment of the literature.
\end{itemize}
\end{document}
